\section*{The Value of Information and Its Spurious Component}

The information laws acquire economic meaning through a simple but exact identity. Consider a price-taking investor with constant absolute risk aversion $\gamma>0$ and initial wealth $w_0$ who can trade $N$ risky assets at prices $p\in\R^N$ and a riskless asset with zero net return. Terminal payoffs are $v\sim\N(\mu,\Sigma_v)$. Before trading, the investor may observe a signal $s\in\R^K$ that is jointly Gaussian with $v$. Let $\Sigma_{v|s}$ denote the conditional covariance of $v$ given $s$, which is non-random, and recall that $I(v;s)=\tfrac12\log\det(\Sigma_v\Sigma_{v|s}^{-1})$. Terminal wealth for a position $x$ is $w_0+x^\top(v-p)$.

\begin{theorem*}[Exact value of information]
Let $U_s$ and $U_\varnothing$ be the ex-ante expected utilities, under $u(w)=-e^{-\gamma w}$, of the optimal policies with and without the signal. Then
\[
 \frac{U_s}{U_\varnothing}=\exp\{-I(v;s)\},
\]
so that the certainty-equivalent value of the signal, the largest sum the investor would pay for it before observing it, is exactly $I(v;s)/\gamma$. The value is independent of $\mu$, $p$ and $w_0$, and is invariant under invertible linear transformations of $s$.
\end{theorem*}

Conditional on $s$, the optimal demand is $x^*=(\gamma\Sigma_{v|s})^{-1}(m-p)$ with $m=\E[v\mid s]$, and the conditional expected utility is $-\exp\{-\gamma w_0-\tfrac12(m-p)^\top\Sigma_{v|s}^{-1}(m-p)\}$. Ex ante, $m-p\sim\N(\mu-p,\Sigma_v-\Sigma_{v|s})$. For $y\sim\N(c,B)$ and positive semidefinite $A$, the Gaussian integral $\E\exp(-\tfrac12y^\top Ay)=\det(I+BA)^{-1/2}\exp\{-\tfrac12c^\top(A^{-1}+B)^{-1}c\}$ holds. With $A=\Sigma_{v|s}^{-1}$ and $B=\Sigma_v-\Sigma_{v|s}$ we have $I+BA=\Sigma_v\Sigma_{v|s}^{-1}$ and $A^{-1}+B=\Sigma_v$. The exponential factor is therefore exactly the uninformed expected utility, and the determinant factor is $e^{-I(v;s)}$. The scalar versions of the two matrix identities are certified in Lean, and Appendix~A gives the details.

The identity says that in the canonical model of portfolio choice under information, mutual information is not merely an analogy for value. It \emph{is} value, measured in units of risk tolerance. It makes precise the sense in which entropy-based measures order signals for investors \citep{cabrales2013entropy}, and it underlies the log-determinant objectives in models of information choice \citep{vannieuwerburgh2010information,veldkamp2011information,kacperczyk2016rational}. It also explains why a linear cost per nat, as in rational inattention \citep{sims2003implications,mackowiak2023rational}, produces bang-bang acquisition decisions in CARA--Gaussian settings: a signal is worth buying at a price of $\lambda$ per nat exactly when $1/\gamma>\lambda$. A further analogy is with \citet{kelly1956new}, who showed that a log-utility gambler with side information grows wealth faster by exactly the mutual information. The theorem is the CARA counterpart, with certainty equivalents in place of growth rates. The identity is partial-equilibrium. In a noisy rational expectations equilibrium prices themselves carry information \citep{grossman1980impossibility,hellwig1980aggregation,verrecchia1982information}, and the relevant quantity becomes the conditional mutual information $I(v;s\mid p)$. The theorem applies verbatim to that quantity because conditioning on a jointly Gaussian price only replaces $\Sigma_v$ by $\Sigma_{v|p}$.

\subsection*{Spurious value of information}

Valuing a signal in practice requires estimating $I(v;s)$. An investor, an econometrician or a data vendor with $n$ joint observations of $(v,s)$, with $\dim v=N$ and $\dim s=K$, who uses the plug-in estimate reports $\hat V=\hat I/\gamma$. All results of the previous section transfer. The expected excess of the reported value over the true one is
\[
 \E\hat V-V=\frac{B_n(N,K)}\gamma>0,
\]
exactly, and the corrected value $\tilde V=(\hat I-B_n(N,K))/\gamma$ is unbiased. Under the null of a worthless signal, $\gamma\hat V$ has the log-Beta law of the null proposition, so a test for positive value is available at any sample size. In the proportional regime the spurious value grows like $n\kappa(N/n,K/n)/\gamma$. It is therefore not a small-sample nuisance that vanishes with more data. If the numbers of assets and signals grow with the sample, it is a permanent wedge, whose size per observation is set by $\kappa$ and shown in Figure~\ref{fig:maps}. With a risk tolerance of one million dollars, the spurious value in the empirical design of the application, with eight assets, four signals and a year of daily data, is $\$\,$\MIBias{}~million, or about $66{,}000$ dollars per decision.

The monotonicity proposition has a sharper implication for markets in which information is sold. Suppose a vendor offers $K$ signals and a buyer evaluates them by the plug-in value on a common sample. Appending $m$ series that are independent of the payoffs and of the original signals leaves the true value unchanged. It raises the expected measured value by
\begin{multline*}
 \Delta_m=\frac{B_n(N,K+m)-B_n(N,K)}\gamma\\
 =\frac1{2\gamma}\sum_{i=1}^m\Bigl[\psi\Bigl(\frac{n-K-i+1}2\Bigr)-\psi\Bigl(\frac{n-N-K-i+1}2\Bigr)\Bigr]>0.
\end{multline*}
A buyer who does not correct for dimension therefore faces a lemons problem in the sense of \citet{akerlof1970market}, in which padding is privately profitable and undetectable on average. The correction restores the right incentives in expectation. It cannot restore precision. The null variance $V_n(N,K+m)$ is strictly increasing in $m$, because $\psi'$ is decreasing, so padded signal sets remain harder to evaluate even when evaluated without bias. This mechanism offers a formal account of a concern that runs through work on data markets and financial data technology \citep{farboodi2020long}: as the number of available series grows, the measured information content of data rises mechanically, while the economically relevant content need not. It is also a precise sense in which uninformative or polluted information sources degrade price discovery. They do not mislead on average once the analyst corrects, but they consume degrees of freedom and inflate uncorrected assessments.

\subsection*{Information that cannot be located}

The frontier \eqref{eq:frontier} adds a further layer. Suppose the dependence between payoffs and signals is concentrated in one direction, so that the true value is $-\tfrac1{2\gamma}\log(1-t^2)$. An investor who knows the direction can trade on it and capture that value. A statistician who must learn the direction from $n$ observations of $N$ payoffs and $K$ signals faces the canonical-correlation phase transition. When $t<t_c$, the leading sample canonical correlation converges to the bulk edge $\omega_+$ whatever the value of $t$, so it reveals nothing about $t$. By analogy with the eigenvector part of the BBP transition, the estimated canonical direction is then expected to carry no asymptotic information about the true direction either. I verify the eigenvalue statement by simulation below, and leave the eigenvector statement as a conjecture supported by the analogy. The value $V_c=I_c(N/n,K/n)/\gamma$ is thus the largest rank-one value that remains inaccessible to learning from data at the given dimensions. Advantages below that level can persist without being competed away by statistical arbitrage. This gives an information-theoretic counterpart to the observation of \citet{martin2022market} that investors who learn in high dimensions face estimation frictions that make returns appear predictable ex post. It also connects to the evidence of \citet{kelly2024virtue} and \citet{kozak2020shrinking} that the relevant signals in asset pricing are numerous and individually weak. In the empirical design used below, the frontier is small, $I_c=\FrontierEmp{}$ nats with $t_c=\FrontierEmpT{}$, because eight and four are small compared with a year of daily data. At the scale of the full cross-section of individual stocks, with hundreds of characteristics, it becomes first order.
